\begin{example}[Identity base change]\label{ex:vi24-r01}
For $X\to S$, pulling back along $\operatorname{id}_S$ gives $X\times_SS\cong X$.
\end{example}

\begin{example}[Affine line over a larger field]\label{ex:vi24-r02}
If $k\subseteq K$, then $\mathbb A^1_k\times_kK\cong\Spec(K[x])=\mathbb A^1_K$.
\end{example}

\begin{example}[Affine space]\label{ex:vi24-r03}
More generally $\mathbb A^n_k\times_kK\cong\mathbb A^n_K$ because $k[x_1,\ldots,x_n]\otimes_kK\cong K[x_1,\ldots,x_n]$.
\end{example}

\begin{example}[Scalar extension of a quotient]\label{ex:vi24-r04}
For $f\in k[x]$, $(k[x]/(f))\otimes_kK\cong K[x]/(f)$; this is the legacy CA13-E7 computation.
\end{example}

\begin{example}[Complexification of a real point]\label{ex:vi24-r05}
$\Spec\mathbb C\times_{\Spec\mathbb R}\Spec\mathbb C\cong\Spec(\mathbb C\times\mathbb C)$, so the base-changed point splits.
\end{example}

\begin{example}[A principal open]\label{ex:vi24-r06}
For $D(b)\subseteq\Spec B$, base change by $A\to A'$ gives $D(b\otimes1)\subseteq\Spec(B\otimes_AA')$.
\end{example}

\begin{example}[A closed subscheme]\label{ex:vi24-r07}
If $Z=\Spec(B/I)\hookrightarrow\Spec B$, then $Z_{A'}\cong\Spec((B\otimes_AA')/I(B\otimes_AA'))$.
\end{example}

\begin{example}[A reducible conic after extension]\label{ex:vi24-r08}
$\Spec\mathbb R[x]/(x^2+1)$ is one closed point over $\mathbb R$, while after base change to $\mathbb C$ the polynomial factors and the scheme becomes two reduced points.
\end{example}

\begin{example}[Dual numbers]\label{ex:vi24-r09}
$k[\varepsilon]/(\varepsilon^2)\otimes_kK\cong K[\varepsilon]/(\varepsilon^2)$, so this nilpotent thickening stays nonreduced after a field extension.
\end{example}

\begin{example}[Base-changing a morphism]\label{ex:vi24-r10}
The $k$-morphism $\mathbb A^1_k\to\mathbb A^1_k$ given by $y=x^2$ becomes the same formula over $K$ after $k\subseteq K$.
\end{example}

\begin{example}[Base change of a product]\label{ex:vi24-r11}
$(\mathbb A^1_k\times_k\mathbb A^1_k)_K\cong\mathbb A^1_K\times_K\mathbb A^1_K\cong\mathbb A^2_K$.
\end{example}

\begin{example}[Base change of the diagonal]\label{ex:vi24-r12}
The equation $x_1-x_2=0$ defining the diagonal of $\mathbb A^1_k$ extends to the same equation in $K[x_1,x_2]$.
\end{example}

\begin{example}[Base change of a graph]\label{ex:vi24-r13}
The graph $y-x^2=0$ in $\mathbb A^2_k$ becomes the graph $y-x^2=0$ in $\mathbb A^2_K$.
\end{example}

\begin{example}[Changing from integers to rationals]\label{ex:vi24-r14}
$\Spec\mathbb Z[x]\times_{\Spec\mathbb Z}\Spec\mathbb Q\cong\Spec\mathbb Q[x]$. This is scalar extension, while its interpretation as the generic fiber is deferred to VI/25.
\end{example}

\begin{example}[Reduction modulo $p$ as scalar extension]\label{ex:vi24-r15}
$\Spec\mathbb Z[x]/(F)\times_{\Spec\mathbb Z}\Spec\mathbb F_p\cong\Spec\mathbb F_p[x]/(\overline F)$. Detailed fiber analysis is VI/25.
\end{example}

\begin{example}[An empty scalar extension]\label{ex:vi24-r16}
If $A' = 0$, then $B\otimes_AA'=0$ and the base-changed affine scheme is empty.
\end{example}

\begin{example}[Disjoint unions]\label{ex:vi24-r17}
Base change commutes with disjoint unions because both are characterized componentwise by universal properties.
\end{example}

\begin{example}[Two equations]\label{ex:vi24-r18}
A scheme $V(f,g)\subseteq\mathbb A^n_k$ base changes to the scheme cut out by the same two polynomials regarded in $K[x_1,\ldots,x_n]$.
\end{example}

\begin{example}[Linear systems]\label{ex:vi24-r19}
A matrix with entries in $k$ defines a linear subscheme; after $k\subseteq K$ the same matrix defines its scalar extension over $K$.
\end{example}

\begin{example}[Base-changing an isomorphism]\label{ex:vi24-r20}
If $X\to Y$ is an isomorphism over $S$, then $X_{S'}\to Y_{S'}$ is an isomorphism with inverse obtained by base-changing the original inverse.
\end{example}
